% Generated draft: requires mathematical and source review.
\documentclass[11pt]{article}
\usepackage[margin=25mm]{geometry}
\usepackage{fontspec}
\setmainfont{Latin Modern Roman}
\usepackage{amsmath,amssymb,mathtools,unicode-math}
\setmathfont{Latin Modern Math}
\usepackage{enumitem,needspace,xurl,hyperref}
\setlength{\parindent}{0pt}
\setlength{\parskip}{5pt}
\setlength{\emergencystretch}{3em}
\newcommand{\N}{\mathbb{N}}
\newcommand{\Z}{\mathbb{Z}}
\newcommand{\Q}{\mathbb{Q}}
\newcommand{\R}{\mathbb{R}}
\newcommand{\C}{\mathbb{C}}
\newcommand{\F}{\mathbb{F}}
\newcommand{\A}{\mathbb{A}}
\newcommand{\PP}{\mathbb{P}}
\newcommand{\E}{\mathbb{E}}
\newcommand{\eps}{\varepsilon}
\newcommand{\dd}{\,\mathrm{d}}
\begin{document}
\section*{UnsolvedMath ID: 30005400}
\textbf{Source problem code:} catalogue entry
\subsection{Definitions and mathematical statement}
A field \(K\) is called \(C_1\) if, for every pair of integers \(d\geq 1\) and \(n>d\), every homogeneous polynomial \(f\in K[T_1,\ldots,T_n]\) of degree \(d\) has a nonzero zero:
\[
\exists\,(a_1,\ldots,a_n)\in K^n\setminus\{(0,\ldots,0)\}
\quad f(a_1,\ldots,a_n)=0.
\]
The classical setting motivating the question is the function field \(k(C)\) of an integral algebraic curve \(C\) over an algebraically closed field \(k\). Tsen's dimension-count argument proves the \(C_1\) property in this setting.

\textbf{Proposed proof-construction problem.} Can Tsen's dimension-count argument be adapted to the Fargues--Fontaine curve by using Banach--Colmez spaces in place of the affine spaces occurring in the classical argument?

To express the intended arithmetic conclusion, let \(X\) denote the intended integral Fargues--Fontaine curve and put \(K=\kappa(\eta_X)\), the residue field of its generic point. The proposed adaptation is intended to establish the displayed \(C_1\) assertion for \(K\). The question concerns constructing such an adaptation, rather than merely establishing the assertion by an unspecified method.

\textbf{Unresolved qualifications.} The supplied record does not specify the defining fields or other base data for \(X\), whether a schematic or analytic version is intended, or the range of choices over which the question is quantified. Thus \(X\) and \(K\) above describe a conditional formulation, not a supplied universal quantification over Fargues--Fontaine curves. The record also does not specify the category of Banach--Colmez spaces, the dimension invariant, or the parameter spaces and vanishing argument needed for the adaptation. Its verbatim excerpt refers to ``this proof'' without reproducing that proof; its identification with Tsen's dimension-count proof comes from the supplied reconstructed statement.
\subsection{Short English statement}
The question asks whether Banach--Colmez spaces can supply a version of Tsen's dimension-count proof showing that every degree-\(d\) homogeneous equation in more than \(d\) variables over the intended Fargues--Fontaine function field has a nonzero solution. The curve's base data and the precise adaptation remain unspecified.
\subsection{Sources}
\begin{itemize}
\item Local registry: \texttt{data/unsolvedmath/problems.json}, permanent UnsolvedMath ID \(30005400\), source problem code \texttt{catalogue entry}, title ``Banach--Colmez Proofs of the Fargues--Fontaine C-One Property.'' The supplied record assigns proposed year \(2023\) and names no individual proposer.
\item Supplied report provenance: \emph{Arithmetic of Shimura Varieties} (2023), Oberwolfach Reports, pp.\ 261--326, DOI \texttt{10.4171/owr/2023/5}. \url{https://doi.org/10.4171/owr/2023/5}. The report documents workshop discussions. No specific page or speaker is supplied for this extracted question.
\item Supplied literature reference: Le Bras, ``Banach--Colmez spaces and coherent sheaves.'' \url{https://arxiv.org/abs/1801.00422}.
\item Supplied literature reference: Colmez--Nizio\l, ``p-adic analytic cohomology.'' \url{https://arxiv.org/abs/2108.12785}.
\item Supplied dated literature curation, checked \texttt{2026-08-22}: the record labels the problem ``partially solved'' and reports that Banach--Colmez spaces are rigorously related to coherent sheaves on the Fargues--Fontaine curve and used in selected pro-\'{e}tale cohomology computations, while no general adaptation of the proposed proof is established. This is supplied dated curation, not an independently verified literature assessment; its status and cited claims require independent verification.
\item Supplied statement-verification metadata labels the reconstruction \texttt{corrected\_verified} and says it was checked against the official Oberwolfach report. This records the supplied provenance claim; no independent source consultation or mathematical verification is asserted here.
\end{itemize}
\end{document}
